\section{Theoretical Support: Implicit Dynamics of In-Context Learning}
\label{appendix:theory_support}

The iterative refinement process in MACER and dual-evolving mechanism is not merely heuristic but possesses theoretical grounding through the lens of implicit in-context learning dynamics. Recent work by \citep{dherin2025learning} demonstrates that transformer-based models can perform in-context learning by implicitly modifying their MLP weights through attention mechanisms. We extend this theoretical framework to explain the convergence properties of our multi-agent reasoning process.

\paragraph{Implicit Weight Updates via Attention Dynamics}
The trajectory history $\mathcal{H}_k$ serves as an \emph{in-context prompt} that induces implicit low-rank updates to the frozen LLM's parameters. Specifically, for a transformer module with MLP layer weights $W$, the context $\mathcal{H}_k$ generates an implicit weight update $\Delta W_k$ through the attention mechanism:

\begin{align}
    \Delta W_k &= \frac{(W \Delta A_k) A(q)^\top}{\|A(q)\|^2}, \nonumber \\
    \text{where} \quad \Delta A_k &= A(\mathcal{H}_k, q) - A(q). \label{eq:delta_W}
\end{align}

Here, $A(\cdot)$ denotes the activation pattern from the attention layer, $A(q)$ represents the baseline activation without context, and $A(\mathcal{H}_k, q)$ captures the contextualized activation with the full reasoning history. The term $\Delta A_k$ quantifies the information injected by the evolving context $\mathcal{H}_k$.
The low-rank nature of $\Delta W_k$ ensures efficient and targeted parameter updates without catastrophic forgetting of pre-trained knowledge.

\paragraph{MDP Policy as an Implicit Function of Context}
Recall from Section~\ref{subsec:macer-process} that the Reflector Agent’s policy $\pi_{\text{ref}}$ maps states $s_k = (q, \mathcal{G}_k, \mathcal{H}_k)$ to actions (sub-queries or \textsc{stop}). Under the implicit learning view, $\pi_{\text{ref}}$ is not a fixed network but an emergent policy $\pi_k$ shaped by $\Delta W_k$. Thus, the sequence $\{\pi_k\}_{k=1}^K$ constitutes a trajectory of implicitly adapted policies driven by the evolving context $\mathcal{H}_k$.

\paragraph{Convergence via Regret Minimization}
We analyze convergence through the lens of episodic regret minimization in the MDP $\mathcal{M} = (\mathcal{S}, \mathcal{A}, P, r)$. Let $V^{\pi}_{s_k} = \mathbb{E}_{\pi}\left[ \sum_{i=k}^{K} \gamma^{i-k} r_i \mid s_k \right]$ denote the value of policy $\pi$ at state $s_k$, and let $V^*_{s_k} = \max_{\pi} V^{\pi}_{s_k}$ be the optimal value. The cumulative regret over $K$ steps is:
\begin{align}
    \mathcal{R}(K) = \sum_{k=1}^{K} \left( V^*_{s_k} - V^{\pi_k}_{s_k} \right).
\end{align}
We establish sublinear regret growth $\mathcal{R}(K) = o(K)$ under the following mild assumptions:

\begin{assumption}[Realizability]
There exists a policy $\pi^*$ such that $\texttt{Suff}(q, \mathcal{G}^*_q) = 1$, and $\pi^*$ is representable by the implicit policy class induced by in-context prompts of the form $(\mathcal{H}; q)$.
\end{assumption}

\begin{assumption}[Bounded Gradient Norm]
The implicit gradient direction $g_k$, defined as the reward-sensitive update signal from $\mathcal{H}_k$, satisfies $\|g_k\| \leq G$ for some constant $G > 0$.
\end{assumption}

Under these assumptions, the following properties hold:

\textbf{Property 1 (Smooth Policy Evolution).}
The value function evolves smoothly with respect to implicit updates:
\begin{align}
    \|V^{\pi_{k+1}} - V^{\pi_k}\|_{\infty} \leq L \|g_k\| + \mathcal{O}(\|g_k\|^2),
\end{align}
for some Lipschitz constant $L > 0$, ensuring stable policy transitions.

\textbf{Property 2 (Expected Policy Improvement).}
Each refinement step yields non-negative expected improvement:
\begin{align}
    \mathbb{E}\left[ V^{\pi_{k+1}}_{s_k} - V^{\pi_k}_{s_k} \mid \mathcal{H}_k \right] \geq \eta \|g_k\|^2 - \sigma_k,
\end{align}
where $\eta > 0$ and $\{\sigma_k\}$ is a martingale difference sequence with $\mathbb{E}[\sigma_k \mid \mathcal{H}_k] = 0$. This follows from the fact that evolving sub-queries generated by the Reflector target knowledge gaps, and the Constructor’s evolving graph refinement increases the likelihood of sufficiency.

\textbf{Property 3 (Vanishing Implicit Gradient).}
As the context becomes increasingly informative, the room for improvement diminishes:
\begin{align}
    \lim_{k \to \infty} \|g_k\| = 0 \quad \text{almost surely}.
\end{align}
This is guaranteed by Assumption 1 (Realizability) and the finite horizon $K$, which ensures the process either reaches a sufficient subgraph ($r_k = 1$) or exhausts its budget.

Together, these properties imply that the sequence $\{\pi_k\}$ converges to a policy $\pi^\dagger$ satisfying $V^{\pi^\dagger}_{s_1} \geq V^*_{s_1} - \epsilon$ for arbitrarily small $\epsilon > 0$ as $K \to \infty$. In practice, with a reasonable horizon (e.g., $K=3$), MACER reliably converges to a sufficient context $\mathcal{G}^*_q$ for faithful answer synthesis.

This analysis establishes that the MACER loop performs an implicit form of policy gradient ascent on the reward landscape defined by context sufficiency, with convergence guarantees rooted in stochastic approximation theory and in-context learning dynamics, providing rigorous foundations for the empirical effectiveness of our reward-based evolving context mechanism.
